\documentclass{zhvt-classic}

% 使用本地源流明体（GenRyuMin2 TW）按文件路径加载，不依赖系统字体名
% 版框（内框／界栏／书口侧栏／鱼尾）改用与版心印章同一印泥棕色：
%   印章图片 yshs.png 实色 = RGB(135,68,52) = HTML 874434
\usepackage{xcolor}
\usepackage{tikz}
% sealbrown 双注册：xcolor 供 TikZ/pgf，L3 供 cls 版框
\definecolor{sealbrown}{HTML}{874434}
\ExplSyntaxOn
\color_set:nnn { sealbrown } { HTML } { 874434 }
\ExplSyntaxOff

\zhvtset{
  grid color       = sealbrown,
  main font        = GenRyuMin2TW-SB,
  main font path   = C:/Users/Administrator/texmf/fonts/opentype/public/GenRyuMin2TW/,
jiazhu font      = GenRyuMin2TW-M,
  jiazhu font path = C:/Users/Administrator/texmf/fonts/opentype/public/GenRyuMin2TW/,
  % 页码框容纳 4 字：正常页「九十七」+前置空 1 格=4 行恰好；
  % 100+ 页「一百零一」共 5 行超出，末字「一/二/三」被框裁掉（2026-10-04 实测），故放宽到 5
  shukou page number chars = 5,
  % 书口书名框高从默认 6 字减到 4 字，把省出的 2 字高拨给 chapter 框，
  % 使 6 字书名+卷名放进 9 字高的 chapter 框后不再压线（2026-10-05）
  shukou title chars = 4,
}

\title{史记}[刺客列传]
\author{汉·司马迁}
\maker{闲叟漫记·重製}


% ===== 索隐述赞样式（2026-11-08 四次定稿 · 愚大人亲手改定参数）=====
% ⚠ 竖排前提：自定义字体家族必须带 RawFeature={vertical:+vert}；
%   simhei.ttf 加 +vhal 会整体加载失败静默回退，只写 +vert；simfang 同理。
\newCJKfontfamily{\shuzanfont}[Path=C:/Windows/Fonts/,RawFeature={vertical:+vert}]{simhei.ttf}
\newCJKfontfamily{\shuzanbody}[Path=C:/Windows/Fonts/,RawFeature={vertical:+vert}]{simfang.ttf}
% 题头：152x30pt 圆角 6pt 棕底白字；baseline (0,2) 使整框落入界行格子不压版框线；
% 白字节点 x 不补偿、y +8.2pt（像素级墨迹实测 dx/dy≈0，L/R 留白约 20pt）。
\newcommand{\shuzanhead}[1]{%
  \begin{tikzpicture}[baseline={(0,2)}]
    \fill[sealbrown, rounded corners=6pt]
      (-76pt,-15pt) rectangle (76pt,15pt);
    \node[text=white, inner sep=0pt, anchor=center] at (0pt,8.2pt)
      {\shuzanfont\fontsize{23}{23}\selectfont #1};
  \end{tikzpicture}%
}
% ★ 修正书口卷名位置：chapter 框起始位置改回 cls 默认值，避免全局影响。
%   只在每个卷首页通过 \juan 命令里的空命令微调该页位置（安全，不会跨页污染）。
\ExplSyntaxOn
\AtBeginDocument{
  \dim_set:Nn \l__zhvt_shukou_chapter_title_start_dim
    { \l__zhvt_shukou_title_height_dim + \l__zhvt_font_size_dim * 2 }
  \dim_set:Nn \l__zhvt_shukou_page_number_start_dim
    { \l__zhvt_shukou_chapter_title_start_dim + \l__zhvt_shukou_chapter_title_height_dim }
}
\ExplSyntaxOff

\begin{document}

% ★★ 版心（书口）排布覆写：把书名放到鱼尾之下（2026-10-04）
%
%   cls 默认版心自上而下：书名（title 静态框，在鱼尾之上）→ 鱼尾 → 篇名（chapter 动态框）→ 页码
%   本模板改为：          鱼尾（位置不动）→ 书名 → 篇名 → 页码分割线 → 页码
%
%   手法：把 cls 放在 title 静态框里的书名，改放进 chapter 动态框（鱼尾之下）；
%         title 静态框留空。书名取 \l__zhvt_title_str，即 \title{...} 的第一参数，
%         所以换书时只改 \title，不必在此重复写书名。
%
%   ⚠ 三条必须遵守的写法（2026-10-04 实测）：
%     ① 汉字必须用 \hbox{...} 包住。flowfram 会把框内容装进一个宽 14.4mm（= 一个 line_sep
%        = 书口栏宽）的纵向 minipage 并按该宽度折行，而整页靠 PDF /Rotate 90 显示、
%        框内文字本身不变向 —— 不包 \hbox 时「围炉夜话」4 字会被逐字折成 4 行，
%        在版面上就成了横排 4 字并横向溢出版框。
%     ② 竖排下**一个 \hbox 就是一整列**。书名与篇名要同处一列，必须放进同一个 \hbox；
%        分成两个 \hbox 时，后一个会被当成新的一列、甩到隔壁栏去
%        （实测「序」跑到书口栏外侧 14.4mm 的正文列上）。中间的 \quad 在框内是横向空白，
%        显示出来正好是书名与篇名之间的纵向间隔。
%     ③ ★ 内容盒要写成 **与框等宽的定宽盒** `\hbox to \l__zhvt_shukou_chapter_title_height_dim{…\hfil}`。
%        实测：chapter 框是把内容盒沿纵向**居中**摆放的，内容宽度一变，书名就跟着漂：
%        自然宽度的 \hbox（书名 38.4mm）落在 x=157.2mm，加篇名后（57.6mm）落在 x=152.4mm，
%        目录页（无篇名）与正文页的书名就差半格。做成等宽盒后，无论有没有篇名，
%        内容固定从框顶起排，实测两页均落在 x=142.8mm，不再漂。
\ExplSyntaxOn
% 把书口 chapter 框（书名+篇名）起始位置下移 1 字高，避免 4 字卷名顶到鱼尾
% ⚠ cls 的 \__zhvt_calculate_dimensions: 在 \AtBeginDocument 里执行，会重算该 dim；
%   必须把覆写挂到 begin document 之后，否则会被二次覆盖（2026-10-05 实测）。
% (旧方案已移到导言区，直接修改 chapter 框起始位置)
% 卷标题命令 #1=书口篇名(鱼尾) #2=正文卷名（含卷序，卷名内空格用 ~）
%
% ★★ 目录页两项修订（2026-10-05，愚大人指出）★★
%   ① 页码起点统一对齐：
%      原状 —— 目录项写成 \par\hspace{4em}卷名\hspace{4em}页码，页码起点 = 卷名实宽 + 4em。
%      卷名实宽随字数变：「卷一~梁惠王上」6 字 +~、「卷七~离娄上」5 字 +~，
%      差 1 个字位 → 卷七~卷十 4 条的页码整体比其余 10 条高 27.21pt（一个字），看着上下不齐。
%      改法 —— 卷名补全角空格（宽 = 1 字 = 27.215pt）补足到 **7 字符位**再写进目录，
%      14 条页码起点遂全部落在同一条竖线上（起点对齐；末字随页码字数自然参差）。
%      ⚠ 只对目录里的副本 \l_tmpc_tl 补，正文卷名仍用 \l_tmpa_tl 原样 —— 正文像素不受影响。
%   ② 卷名整体上移半字位：前导 \hspace 由 4em 改 3.5em。
%      实测原状目录块在版框内偏下（天头留白 139.6pt / 地脚留白 110.7pt），
%      上移 0.5em = 13.6pt 后上下基本均分，视觉居中。
\NewDocumentCommand { \juan } { mm } {
	\clearpage
	\tl_set:Nn \l__zhvt_chapter_title_tl { #1 }
	\tl_set:Nn \l_tmpa_tl { #2 }
	\tl_set_eq:NN \l_tmpc_tl \l_tmpa_tl
	\int_step_inline:nn { 7 - \str_count:n { #2 } }
	  { \tl_put_right:Nn \l_tmpc_tl { 　 } }
	\tl_set:Nx \l_tmpb_tl { \zhnum { page } }
	\tl_put_left:Nn \l_tmpb_tl { \hspace { 4em } }
	\tl_put_left:Nn \l_tmpb_tl { \l_tmpc_tl }
	\tl_put_left:Nn \l_tmpb_tl { \par \hspace { 3.5em } }
	\addtocontents { toc } { \l_tmpb_tl }
	\skip_horizontal:n { \l__zhvt_font_size_dim * 4 }
	\l_tmpa_tl
	\par
	\thispagestyle { empty }
}
\cs_set_protected:Npn \__zhvt_set_shukou_content: {
	% ★ 鱼尾上方留白（title 静态框）嵌入闲叟钤印：
	%   整页 /Rotate 90，故此处把图用 \rotatebox{90} 转正为竖章；
	%   尺寸卡在版心条幅内（书内约 7mm 宽 × 33.8mm 高，留余量不压鱼尾）。
	\setstaticcontents * { title } {
		\vbox_to_ht:nn { \l__zhvt_baseline_skip_dim } {
			\vskip 14.2pt
			\hbox_to_wd:nn { \l__zhvt_shukou_title_height_dim } {
				\hfil
				\rotatebox [ origin = c ] { 90 } { \includegraphics [ width = 7mm , height = 33.8mm ] { yshs.png } }
				\hfil
			}
		}
	}
% ★ 版心（书口）统一改法（2026-11-16 愚大人定）：18pt＋框内居中＋书口文字同版框色
%   ① 原 27.2pt 时 7 字长篇名（如「廉颇蔺相如列传」）末字溢出 chapter 框、压住页码；
%      改 \fontsize{18pt}{18pt}，前后 \hfil 令字串在框内居中（不随篇名字数漂移）。
%   ② \vbox{\vskip 8.2pt} 定量下移，抵消缩字号后栏宽方向 8.2pt 上偏（像素实测 y 中心 543.3 = 原版）。
%      ⚠ \raisebox / \vspace* 在 flowfram 框内无效；换字号须同步换补偿：16pt→10.1、19pt→7.4。
%   ③ chapter 与 page~number 均置 \color{sealbrown}（HTML 874434），与版框／界行同色；正文仍黑。
%   ⚠ 页码字位、钤印、鱼尾、界行、正文一律未动；五部合订本同此口径。
	\setdynamiccontents * { chapter }
	{ \color{sealbrown} \vbox { \vskip 8.2pt \hbox to \l__zhvt_shukou_chapter_title_height_dim
		{ \hfil \fontsize{18pt}{18pt}\selectfont \l__zhvt_title_str \quad \l__zhvt_chapter_title_tl \hfil } } }
	\setdynamiccontents * { page~number } { \color{sealbrown} \quad \zhnum { page } }
}
\__zhvt_set_shukou_content:
\ExplSyntaxOff

\maketitle

% 单篇，无多卷目录，不渲染目录页

\juan{刺客列传}{卷八十六~刺客列传第二十六}
\hspace*{2em}曹沫者。鲁人也。以勇力事鲁庄公。庄公好力。曹沫为鲁将。与齐战。三败北。鲁庄公惧。乃献遂邑之地以和。犹复以为将。\\
\hspace*{2em}齐桓公许与鲁会于柯而盟。桓公与庄公既盟于坛上。曹沫执匕首劫齐桓公。桓公左右莫敢动。而问曰。﹁子将何欲。﹂曹沫曰。﹁齐强鲁弱。而大国侵鲁亦甚矣。今鲁城坏即压齐境。君其图之。﹂桓公乃许尽归鲁之侵地。既已言。曹沫投其匕首。下坛．北面就群臣之位．颜色不变。辞令如故。桓公怒。欲倍其约。管仲曰。﹁不可。夫贪小利以自快。弃信于诸侯。失天下之援。不如与之。﹂于是桓公乃遂割鲁侵地。曹沫三战所亡地尽复予鲁。\\
\hspace*{2em}其后百六十有七年而吴有专诸之事。\\
\hspace*{2em}专诸者。吴堂邑人也。伍子胥之亡楚而如吴也。知专诸之能。伍子胥既见吴王僚。说以伐楚之利。吴公子光曰。﹁彼伍员父兄皆死于楚而员言伐楚。欲自为报私雠也。非能为吴。﹂吴王乃止。伍子胥知公子光之欲杀吴王僚。乃曰。﹁彼光将有内志。未可说以外事。﹂乃进专诸于公子光。\\
\hspace*{2em}光之父曰吴王诸樊。诸樊弟三人。次曰余祭。次曰夷昧。次曰季子札。诸樊知季子札贤而不立太子。以次传三弟。欲卒致国于季子札。诸樊既死。传余祭。余祭死。传夷昧。夷昧死。当传季子札。季子札逃不肯立。吴人乃立夷昧之子僚为王。公子光曰。﹁使以兄弟次邪。季子当立。必以子乎。则光真适嗣。当立。﹂故尝阴养谋臣以求立。\\
\hspace*{2em}光既得专诸。善客待之。九年而楚平王死。春。吴王僚欲因楚丧。使其二弟公子盖余。属庸将兵围楚之灊。使延陵季子于晋。以观诸侯之变。楚发兵绝吴将盖余。属庸路。吴兵不得还。于是公子光谓专诸曰。﹁此时不可失。不求何获!且光真王嗣。当立。季子虽来。不吾废也。﹂专诸曰。﹁王僚可杀也。母老子弱。而两弟将兵伐楚。楚绝其后。方今吴外困于楚。而内空无骨鲠之臣。是无如我何。﹂公子光顿首曰。﹁光之身。子之身也。﹂\\
\hspace*{2em}四月丙子。光伏甲士于窟室中。而具酒请王僚。王僚使兵陈自宫至光之家。门户阶陛左右。皆王僚之亲戚也。夹立侍。皆持长铍。酒既酣。公子光详为足疾。入窟室中。使专诸置匕首鱼炙之腹中而进之。既至王前。专诸擘鱼。因以匕首刺王僚。王僚立死。左右亦杀专诸。王人扰乱。公子光出其伏甲以攻王僚之徒。尽灭之。遂自立为王。是为阖闾。阖闾乃封专诸之子以为上卿。\\
\hspace*{2em}其后七十余年而晋有豫让之事。\\
\hspace*{2em}豫让者。晋人也。故尝事范氏及中行氏。而无所知名．去而事智伯。智伯甚尊宠之。及智伯伐赵襄子。赵襄子与韩。魏合谋灭智伯。灭智伯之后而三分其地。赵襄子最怨智伯。漆其头以为饮器。豫让遁逃山中。曰。﹁嗟乎!士为知己者死。女为说己者容。今智伯知我。我必为报仇而死。以报智伯。则吾魂魄不愧矣。﹂乃变名姓为刑人。入宫涂厕。中挟匕首。欲以刺襄子。襄子如厕。心动。执问涂厕之刑人。则豫让。内持刀兵。曰。﹁欲为智伯报仇。﹂左右欲诛之。襄子曰。﹁彼义人也。吾谨避之耳。且智伯亡无后。而其臣欲为报仇。此天下之贤人也。﹂卒释去之。\\
\hspace*{2em}居顷之。豫让又漆身为厉。吞炭为哑。使形状不可知。行乞于市。其妻不识也。行见其友。其友识之。曰。﹁汝非豫让邪。﹂曰。﹁我是也。﹂其友为泣曰。﹁以子之才。委质而臣事襄子。襄子必近幸子。近幸子。乃为所欲。顾不易邪。何乃残身苦形。欲以求报襄子。不亦难乎。﹂豫让曰。﹁既已委质臣事人。而求杀之。是怀二心以事其君也。且吾所为者极难耳。然所以为此者。将以愧天下后世之为人臣怀二心以事其君者也。﹂\\
\hspace*{2em}既去。顷之。襄子当出。豫让伏于所当过之桥下。襄子至桥。马惊。襄子曰。﹁此必是豫让也。﹂使人问之。果豫让也。于是襄子乃数豫让曰。﹁子不尝事范。中行氏乎。智伯尽灭之。而子不为报仇。而反委质臣于智伯。智伯亦已死矣。而子独何以为之报仇之深也。﹂豫让曰。﹁臣事范。中行氏。范。中行氏皆众人遇我。我故众人报之。至于智伯。国士遇我。我故国士报之。﹂襄子喟然叹息而泣曰。﹁嗟乎豫子。子之为智伯。名既成矣。而寡人赦子。亦已足矣。子其自为计。寡人不复释子。﹂使兵围之。豫让曰。﹁臣闻明主不掩人之美。而忠臣有死名之义。前君已宽赦臣。天下莫不称君之贤。今日之事。臣固伏诛。然愿请君之衣而击之．焉以致报仇之意。则虽死不恨。非所敢望也。敢布腹心。﹂于是襄子大义之。乃使使持衣与豫让。豫让拔剑三跃而击之。曰。﹁吾可以下报智伯矣。﹂遂伏剑自杀。死之日。赵国志士闻之。皆为涕泣。\\
\hspace*{2em}其后四十余年而轵有聂政之事。\\
\hspace*{2em}聂政者。轵深井里人也。杀人避仇。与母。姊如齐。以屠为事。\\
\hspace*{2em}久之。濮阳严仲子事韩哀侯。与韩相侠累有郤。严仲子恐诛。亡去。游求人可以报侠累者。至齐。齐人或言聂政勇敢士也。避仇隐于屠者之间。严仲子至门请。数反。然后具酒自畅聂政母前。酒酣。严仲子奉黄金百溢。前为聂政母寿。聂政惊怪其厚。固谢严仲子。．严仲子固进。而聂政谢曰。﹁臣幸有老母。家贫。客游以为狗屠。可以旦夕得甘毳以养亲。亲供养备。不敢当仲子之赐。﹂严仲子辟人。因为聂政言曰。﹁臣有仇。而行游诸侯众矣。然至齐。窃闻足下义甚高。故进百金者。将用为大人粗粝之费。得以交足下之驩。岂敢以有求望邪。﹂聂政曰。﹁臣所以降志辱身居市井屠者．徒幸以养老母。老母在。政身未敢以许人也。﹂严仲子固让。聂政竞不肯受也。然严仲子卒备宾主之礼而去。\\
\hspace*{2em}久之。聂政母死。既已葬。除服。聂政曰。﹁嗟乎!政乃市井之人。鼓刀以屠。而严仲子乃诸侯之卿相也。不远千里。枉车骑而交臣。臣之所以待之。至浅鲜矣。未有大功可以称者。而严仲子奉百金为亲寿。我虽不受。然是者徒深知政也。夫贤者以感忿睚眦之意而亲信穷僻之人。而政独安得嘿然而已乎。且前日要政。政徒以老母。老母今以天年终。政将为知己者用。﹂乃遂西至濮阳。见严仲子曰。﹁前日所以不许仲子者。徒以亲在。今不幸而母以天年终。仲子所欲报仇者为谁。请得从事焉。﹂严仲子具告曰。﹁臣之仇韩相侠累。侠累又韩君之季父也。宗族盛多。居处兵卫甚设。臣欲使人刺之。终莫能就。今足下幸而不弃。请益其车骑壮士可为足下辅翼者。﹂聂政曰。﹁韩之与卫。相去中间不甚远。今杀人之相。相又国君之亲。此其势不可以多人。多人不能无生得失。生得失则语泄。语泄是韩举国而与仲子为雠。岂不殆哉。﹂遂谢车骑人徒。聂政乃辞独行。\\
\hspace*{2em}杖剑至韩。韩相侠累方坐府上。持兵戟而卫侍者甚众。聂政直入。上阶刺杀侠累。左右大乱聂政大呼。所击杀者数十人。因自皮面决眼。自屠出肠。遂以死。\\
\hspace*{2em}韩取聂政尸暴于市。购问莫知谁子。于是韩县购之。有能言杀相侠累者予千金。久之莫知也。\\
\hspace*{2em}政姊荣闻人有刺杀韩相者。贼不得。国不知其名姓。暴其尸而县之千金。乃於邑曰。﹁其是吾弟与。嗟乎。严仲子知吾弟!﹂立起。如韩。之市。而死者果政也。伏尸哭极哀。曰。﹁是轵深井里所谓聂政者也。﹂市行者诸众人皆曰。﹁此人暴虐吾国相。王县购其名姓千金。夫人不闻与。何敢来识之也。﹂荣应之曰。﹁闻之。然政所以蒙污辱自弃于市贩之间者。为老母幸无恙。妾未嫁也。亲既以天年下世。妾已嫁夫。严仲子乃察举吾弟困污之中而交之。泽厚矣。可奈何!士固为知己者死。今乃以妾尚在之故。重自刑以绝从。妾其奈何畏殁身之诛。终灭贤弟之名。﹂大惊韩市人。乃大呼天者三。卒於邑悲哀而死政之旁。\\
\hspace*{2em}晋。楚。齐。卫闻之。皆曰。﹁非独政能也．乃其姊亦烈女也。乡使政诚知其姊无濡忍之志。不重暴骸之难。必绝险千里以列其名。姊弟俱谬于韩市者。亦未必敢以身许严仲子也。严仲子亦可谓知人能得士矣。﹂\\
\hspace*{2em}其后二百二十余年秦有荆轲之事。\\
\hspace*{2em}荆轲者。卫人也。其先乃齐人。徙于卫。卫人谓之庆卿。而之燕。燕人谓之荆卿。荆卿好读书击剑。以术说卫元君。卫元君不用。其后秦伐魏。置东郡。徙卫元君之支属于野王。\\
\hspace*{2em}荆轲尝游过榆次。与盖聂论剑。盖聂怒而目之。荆轲出。人或言复召荆卿。盖曰。﹁曩者吾与论剑有不称者。吾目之。试往。是宜去。不敢留。﹂使使往之主人。荆卿则已驾而去榆次矣。使者还报。盖聂曰。﹁固去也。吾曩者目摄之。﹂\\
\hspace*{2em}荆轲游于邯郸。鲁勾践与荆轲博。争道。鲁勾践怒而叱之。荆轲嘿而逃去。遂不复会。\\
\hspace*{2em}荆轲既至燕。爱燕之狗屠及善击筑者高渐离。荆轲嗜酒。日与狗屠及高渐离饮于燕市。酒酣以往。高渐离击筑。荆轲和而歌于市中。相乐也。已而相泣。旁若无人者。荆轲虽游于酒人乎。然其为人沉深好书。其所游诸侯。尽与其贤豪长者相结。其之燕．燕之处士田光先生亦善待之。知其非庸人也。\\
\hspace*{2em}居顷之。会燕太子丹质秦亡归燕。燕太子丹者。故尝质于赵。而秦王政生于赵。其少时与丹欢。及政立为秦王。而丹质于秦。秦王之遇燕太子丹不善。故丹怨而亡归。\\
\hspace*{2em}归而求为报秦王者。国小。力不能。其后秦日出兵山东以伐齐。楚。三晋。稍蠹食诸侯。且至于燕。燕君臣皆恐祸之至。太子丹患之。问其傅鞠武。武对曰。秦地遍天下。威胁韩。魏。赵氏。北有甘泉。谷口之固。南有泾。渭之沃。擅巴。汉之饶。右陇。蜀之山。左关。殽之险。民众而士厉。兵革有余。意有所出。则长城之南。易水以北。未有所定也。奈何以见陵之怨。欲批其逆鳞哉。丹曰。﹁然则何由。﹂对曰。﹁请入图之。﹂\\
\hspace*{2em}居有间。秦将樊於期得罪于秦王。亡之燕。太子受而舍之。鞠武谏曰。﹁不可。夫以秦王之暴而积怒于燕。足为寒心。又况闻樊将军之所在乎。是谓﹃委肉当饿虎之蹊﹄也。祸必不振矣。虽有管。晏。不能为之谋也。愿太子疾遣樊将军入匈奴以灭口。请西约三晋。南连齐。楚。北购于单于。其后乃可图也。﹂太子曰。﹁太傅之计。旷日弥久。心惛然。恐不能须臾。且非独于此也。夫樊将军穷困于天下。归身于丹。丹终不以迫于强秦而弃所哀怜之交。置之匈奴。是固丹命卒之时也。愿太傅更虑之。﹂鞠武曰。﹁夫行危欲求安。造祸而求福。计浅而怨深。连结一人之后交。不顾国家之大害。此所谓﹃资怨而助祸﹄矣。夫以鸿毛燎于炉炭之上。必无事矣。且以雕鸷之秦。行怨暴之怒。岂足道哉。燕有田光先生。其为人智深而勇沉。可与谋。﹂太子曰。﹁愿因太傅而得交于田先生。可乎。﹂鞠武曰。﹁敬诺。﹂出见田先生。道﹁太子愿图国事于先生也﹂。田光曰。﹁敬奉教。﹂乃造焉。\\
\hspace*{2em}太子逢迎。却行为导。跪而蔽席。田光坐定。左右无人。太子避席而请曰。﹁燕秦不两立。愿先生留意也。﹂田光曰。﹁臣闻骐骥盛壮之时。一日而驰千里。至其衰老。驽马先之。今太子闻光盛壮之时。不知臣精已消亡矣。虽然。光不敢以图国事。所善荆卿可使也。﹂太子曰。﹁愿因先生得结交于荆卿。可乎。﹂田光曰。﹁敬诺。﹂即起。趋出。太子送至门。戒曰。﹁丹所报。先生所言者。国之大事也。愿先生勿泄也。﹂田光俛而笑曰。﹁诺。﹂偻行见荆卿。曰。﹁光与子相善。燕国莫不知。今太子闻光壮盛之时。不知吾形已不逮也。幸而教之曰﹃燕秦不两立。愿先生留意也﹄。光窃不自外。言足下于太子也。愿足下过太子于宫。﹂荆轲曰。﹁谨奉教。﹂田光曰。﹁吾闻之。长者为行。不使人疑之。今太子告光曰﹃所言者。国之大事也。愿先生勿泄﹄。是太子疑光也。夫为行而使人疑之。非节侠也。﹂欲自杀以激荆卿。曰。﹁愿足下急过太子。言光已死。明不言也。﹂因遂自刎而死。\\
\hspace*{2em}荆轲遂见太子。言田光已死。致光之言。太子再拜而跪。膝行流涕。有顷而后言曰。﹁丹所以诫田先生毋言者。欲以成大事之谋也。今田先生以死明不言。岂丹之心哉!﹂荆轲坐定。太子避席顿首曰。﹁田先生不知丹之不肖。使得至前。敢有所道。此天之所以哀燕而不弃其孤也。今秦有贪利之心。而欲不可足也。非尽天下之地。臣海内之王者。其意不厌。今秦已虏韩王。尽纳其地。又举兵南伐楚。北临赵。王翦将数十万之众距漳。邺。而李信出太原。云中。赵不能支秦。必入臣。入臣则祸至燕。燕小弱。数困于兵。今计举国不足以当秦。诸侯服秦。莫敢合从。丹之私计愚。以为诚得天下之勇士使于秦。窥以重利。秦王贪。其势必得所愿矣。诚得劫秦王。使悉反诸侯侵地。若曹沫之与齐桓公。则大善矣。则不可。因而刺杀之。彼秦大将擅兵于外而内有乱。则君臣相疑。以其间诸侯得合从。其破秦必矣。此丹之上愿。而不知所委命。唯荆卿留意焉。﹂\\
\hspace*{2em}久之。荆轲曰。﹁此国之大事也。臣驽下。恐不足任使。﹂太子前顿首。固请毋让。然后许诺。于是尊荆卿为上卿。舍上舍。太子日造门下。供太牢具。异物间进。车骑美女恣荆轲所欲。以顺适其意。久之。荆轲未有行意。秦将王翦破赵。虏赵王。尽收入其地。进兵北略地至燕南界。太子丹恐惧。乃请荆轲曰。﹁秦兵旦暮渡易水。则虽欲长侍足下。岂可得哉!﹂荆轲曰。﹁微太子言。臣愿谒之。今行而毋信。则秦未可亲也。夫樊将军。秦王购之金千斤。邑万家。诚得樊将军首与燕督亢之地图。奉献秦王。秦王必说见臣。臣乃得有以报。﹂太子曰。﹁樊将军穷困来归丹。丹不忍以己之私而伤长者之意。愿足下更虑之。﹂\\
\hspace*{2em}荆轲知太子不忍。乃遂私见樊於期曰。﹁秦之遇将军可谓深矣。父母宗族皆为戮没。今闻购将军首金千斤。邑万家。将奈何?﹂於期仰天太息流涕曰。﹁於期每念之。常痛于骨髓。顾计不知所出耳。﹂荆轲曰。﹁今有一言可以解燕国之患。报将军之仇者。何如?﹂於期乃前曰。﹁为之奈何。﹂荆轲曰。﹁愿得将军之首以献秦王。秦王必喜而见臣。臣左手把其袖。右手揕其匈。然则将军之仇报而燕见陵之愧除矣。将军岂有意乎。﹂樊於期偏袒扼腕而进曰。﹁此臣之日夜切齿腐心也。乃今得闻教!﹂遂自刭。太子闻之。驰往。伏尸而哭。极哀。既已不可奈何。乃遂盛樊於期首函封之。\\
\hspace*{2em}于是太子豫求天下之利匕首。得赵人徐夫人匕首。取之百金。使工以药蟀之。以试人。血濡缕。人无不立死者。乃装为遣荆卿。燕国有勇士秦武阳。年十三。杀人。人不敢忤视。乃令秦武阳为副。荆轲有所待。欲与俱。其人居远未来。而为治行。顷之。未发。太子迟之。疑其改悔。乃复请曰。﹁曰已尽矣。荆卿岂有意哉。丹请得先遣秦武阳。﹂荆轲怒。叱太子曰。﹁何太子之遣。往而不返者。竖子也。且提一匕首入不测之强秦。仆所以留者。待吾客与俱。今太子迟之。请辞决矣。﹂遂发。\\
\hspace*{2em}太子及宾客知其事者。皆白衣冠以送之。至易水之上。既祖。取道。高渐离击筑。荆轲和而歌。为变徵之声。士皆垂泪涕泣。又前而为歌曰。﹁风萧萧兮易水寒。壮士一去兮不复还。﹂复为羽声伉慨。士皆瞋目。发尽指冠。于是荆轲就车而去。终已不顾。\\
\hspace*{2em}遂至秦。持千金之资币物。厚遗秦王宠臣中庶子蒙嘉。嘉为先言于秦王曰。﹁燕王诚振怖大王之威。不敢举兵以逆军吏。愿举国为内臣。比诸侯之列。给贡职如郡县。而得奉守先王之宗庙。恐惧不敢自陈。谨斩樊於期之头。及献燕督亢之地图．函封。燕王拜送于庭。使使以闻大王。唯大王命之。﹂秦王闻之。大喜。乃朝服。设九宾。见燕使者咸阳宫。荆轲奉樊於期头函。而秦舞阳奉地图柙。以次进。至陛。秦武阳色变振恐。群臣怪之。荆轲顾笑武阳。前谢曰。﹁北蕃蛮夷之鄙人。未尝见天子。故振慴。愿大王少假借之。使得毕使于前。﹂秦王谓轲曰。﹁取武阳所持地图。﹂轲既取图奏之。秦王发图。图穷而匕首见。因左手把秦王之袖。而右手持匕首椹之。未至身。秦王惊。自引而起。袖绝。拔剑。剑长。操其室。时惶急。剑坚。故不可立拔。荆轲逐秦王。秦王环柱而走。群臣皆愕。卒起不意。尽失其度。而秦法。群臣侍殿上者不得持尺寸之兵。诸郎中执兵皆陈殿下。非有诏召不得上。方急时。不及召下兵。以故荆轲乃逐秦王。而卒惶急。无以击轲。而以手共搏之。是时侍医夏无且以其所奉药囊提荆轲也。秦王方环柱走。卒惶急。不知所为。左右乃曰。﹁王负剑。﹂负剑。遂拔以击荆轲．断其左股。荆轲废。乃引其首以擿秦王。不中。中桐柱。秦王复击轲。轲被八创。轲自知事不就。倚柱而笑．箕踞以骂曰。﹁事所以不成者。以欲生劫之。必得约契以报太子也。﹂于是左右既前杀轲。秦王不怡者良久。已而论功。赏群臣及当坐者各有差。而赐夏无且黄金二百溢。曰。﹁无且爱我。乃才以药囊提荆轲也。﹂\\
\hspace*{2em}于是秦王大怒。益发兵诣赵。诏王翦军以伐燕。十月而拔蓟城。燕王喜。太子丹等尽率其精兵东保于辽东。秦将李信追击燕王急。代王嘉乃遗燕王喜书曰。﹁秦所以尤追燕急者。以太子丹故也。今王诚杀丹献之秦王。秦王必解。而社稷幸得血食。﹂其后李信追丹。丹匿衍水中。燕王乃使使斩太子丹。欲献之秦。秦复进兵攻之。后五年。秦卒灭燕。虏燕王喜。\\
\hspace*{2em}其明年。秦并天下。立号为皇帝。于是秦逐太子丹。荆轲之客。皆亡。高渐离变名姓为人庸保。匿作于宋子。久之。作苦。闻其家堂上客击筑。傍惶不能去。每出言曰。﹁彼有善有不善。﹂从者以告其主。曰。﹁彼庸乃知音。窃言是非。﹂家丈人召使前击筑。一坐称善。赐酒。而高渐离念久隐畏约无穷时。乃退。出其装匣中筑与其善衣。更容貌而前。举坐客皆惊。下与抗礼。以为上客。使击筑而歌。客无不流涕而去者。宋子传客之。闻于秦始皇。秦始皇召见。人有识者。乃曰。﹁高渐离也。﹂秦皇帝惜其善击筑。重赦之。乃嚯其目。使击筑。未尝不称善。稍益近之。高渐离乃以铅置筑中。复进得近。举筑朴秦皇帝。不中。于是遂诛高渐离。终身不复近诸侯之人。\\
\hspace*{2em}鲁勾践已闻荆轲之刺秦王。私曰。﹁嗟乎。惜哉其不讲于刺剑之术也!甚矣吾不知人也!曩者吾叱之。彼乃以我为非人也。﹂\\
\hspace*{2em}太史公曰。世言荆轲。其称太子丹之命。﹁天雨粟。马生角﹂也。太过。又言荆轲伤秦王。皆非也。始公孙季功。董生与夏无且游。具知其事。为余道之如是。自曹沫至荆轲五人。此其义或成或不成。然其立意较然。不欺其志。名垂后世。岂妄也哉。\\

% ===== 索隐述赞（司马贞《史记索隐》）· 统一样式 =====
% 题头「索隐述赞曰」五字并入反白框；正文每四句一列，定宽盒 \hfil 居列中；
% \fontsize{24}{41} 列距对齐全书界行（121+41k）；\clearpage 使述赞独占一页。
\clearpage
\shuzanhead{索隐述赞曰}\linebreak
{\shuzanbody\fontsize{24}{41}\selectfont
\hbox to \hsize{\hfil 曹沫盟柯。返鲁侵地。专诸进炙。定吴篡位。\hfil}\linebreak
\hbox to \hsize{\hfil 彰弟哭市。报主涂厕。刎颈申冤。操袖行事。\hfil}\linebreak
\hbox to \hsize{\hfil 暴秦夺魄。懦夫增气。\hfil}%
\par}
\end{document}
